\sectioncentered*{Введение}
\addcontentsline{toc}{section}{Введение}

Продажа подержанных товаров в Беларуси в значительной мере переместилась на площадки частных объявлений. Крупнейшая из них~-- \textit{Kufar.by}, где одновременно размещены десятки тысяч объявлений в каждой из популярных категорий. Покупатель, который ищет товар дешевле рыночной цены, вынужден по нескольку раз в день вручную просматривать выдачу. Новые объявления появляются круглосуточно, а предложения с заниженной ценой нередко снимаются с публикации в течение нескольких часов. Кто открыл площадку позже, узнаёт о таком объявлении уже после сделки. Отсюда потребность в системе, которая следит за площадкой вместо человека и сообщает ему только о подходящих предложениях.

Внешне такая система сводится к интерфейсу: пользователь задаёт в \textit{Telegram} товар и пороговую цену, бот присылает ссылки на объявления. Но на интерфейс приходится меньшая часть работы. Основная нагрузка связана с обработкой данных, которые поступают без перерыва. Площадка опрашивается каждые несколько минут, и каждый полученный ответ сопоставляется с уже накопленными сведениями. Система должна отличить новое объявление от просмотренного ранее, заметить снижение цены на знакомое объявление, помнить, по какой цене предложение уже отправлялось пользователю, и не присылать одно и то же уведомление дважды. Каждая из этих операций сводится к запросу к хранимым данным, поэтому по своей сути задача относится к управлению данными.

Файловое хранение для этих данных не подходит. С ними одновременно работают несколько процессов: фоновая задача опроса записывает объявления и найденные совпадения, задача доставки отмечает отправленные уведомления, а сервис, обслуживающий бота, читает и изменяет настройки пользователей. Без согласованного параллельного доступа эти операции мешали бы друг другу. Записи связаны в цепочку из пользователя, его отслеживаемых товаров, совпадений с объявлениями и отправленных уведомлений, и удаление товара не должно оставлять в хранилище совпадения, которые ссылаются на несуществующую запись. Доставка уведомления требует транзакции: отметка об отправке и запись в журнал уведомлений фиксируются вместе, иначе сбой между двумя шагами приведёт к потере сообщения или к его повтору. Параллельный доступ, ссылочную целостность и транзакции обеспечивает СУБД, поэтому основой системы служит реляционная база данных \textit{PostgreSQL}.

Цель курсового проекта~-- разработка автоматизированной информационной системы отслеживания цен на товары в виде \textit{Telegram}-бота, который ведёт мониторинг объявлений \textit{Kufar.by} и уведомляет пользователей о предложениях ниже заданной цены.

Для достижения цели поставлены следующие задачи:
\begin{itemize}
	\item проанализировать предметную область и существующие средства мониторинга объявлений;
	\item спроектировать базу данных и архитектуру системы;
	\item реализовать базу данных, периодический опрос площадки, отбор подходящих объявлений и доставку уведомлений;
	\item подготовить руководство для пользователя и администратора.
\end{itemize}
